\chapter{Introdução}
\label{sec:introducao}

Na era digital contemporânea, a produção e o consumo de conteúdo de mídia, especialmente vídeos e áudios, atingiram uma escala sem precedentes. Plataformas de educação, entretenimento e comunicação corporativa dependem cada vez mais da capacidade de extrair informações valiosas e insights de grandes volumes de dados não estruturados. Nesse contexto, a automação de tarefas como transcrição e sumarização de conteúdo tornou-se não apenas uma conveniência, mas uma necessidade estratégica. No entanto, a construção de sistemas capazes de realizar essas tarefas de forma eficiente, escalável e, acima de tudo, flexível, apresenta desafios arquiteturais significativos.

Sistemas de processamento de mídia tradicionais frequentemente adotam arquiteturas monolíticas e fortemente acopladas, onde o fluxo de processamento é rígido e focado em um único tipo de mídia, geralmente o vídeo. Essa abordagem resulta em sistemas frágeis e de difícil manutenção, onde a adição de um novo tipo de mídia (como um arquivo de áudio ou um documento PDF) ou a integração de um novo serviço de IA (como um provedor de transcrição alternativo) exige refatorações complexas e arriscadas. A falta de separação de responsabilidades e o acoplamento direto entre os componentes do pipeline limitam a escalabilidade e a resiliência do sistema como um todo.

Este trabalho aborda esses desafios através do design e da implementação de uma plataforma de processamento de mídia genérica, fundamentada em uma arquitetura modular, orientada a eventos e baseada em padrões de design de software modernos. A solução proposta abstrai o conceito de "vídeo" para uma entidade de "Mídia" mais genérica, permitindo que o sistema orquestre pipelines de processamento dinâmicos e desacoplados para diferentes formatos de entrada. A comunicação entre os módulos especializados (e.g., extração de áudio, transcrição, sumarização) é realizada exclusivamente através de um barramento de eventos, garantindo alta coesão e baixo acoplamento.

A relevância deste projeto reside na aplicação prática e na síntese de um conjunto de padrões de arquitetura avançados, como Command Query Responsibility Segregation (CQRS), o Padrão State para gerenciamento de ciclo de vida, e o Padrão Strategy para a seleção dinâmica de provedores de serviço. O resultado é uma plataforma que não apenas cumpre os requisitos funcionais de transcrição e sumarização, mas o faz de uma maneira que é inerentemente extensível, resiliente e manutenível, servindo como um estudo de caso para a construção de sistemas distribuídos modernos.

Este documento está estruturado da seguinte forma: o Capítulo \ref{sec:objetivos} detalha os objetivos geral e específicos do projeto. O Capítulo \ref{sec:definicoes} apresenta o problema, a proposta de solução e as tecnologias utilizadas. O Capítulo \ref{sec:requisitos} descreve os requisitos funcionais e não funcionais. O Capítulo \ref{sec:design} aprofunda-se no design do sistema, detalhando a arquitetura de software e o modelo de banco de dados. Os capítulos subsequentes abordam a qualidade, a implantação e o manual do usuário, concluindo com as considerações finais e os resultados alcançados.
